\subsection{PROPERTIES OF UNION AND INTERSECTION}

If \((\mathbf{X}_t)_{t\in \mathbf{I}}\) and \((\mathbf{Y}_t)_{t\in \mathbf{I}}\) are families of sets having the same index set \(\mathbf{I}\) , and if \(\mathbf{Y}_t\subset \mathbf{X}_t\) for each \(t\in \mathbf{I}\) , then it is clear that

\[
\bigcup_ {i \in I} Y _ {i} \subset \bigcup_ {i \in I} X _ {i}, \quad \text { and } (\text { if   } I \neq \emptyset) \quad \bigcap_ {i \in I} Y _ {i} \subset \bigcap_ {i \in I} X _ {i}.
\]

Let \((\mathbf{X}_t)_{t\in \mathbf{I}}\) be a family of sets. If \(J\subset I\) , we have

\[
\bigcup_ {i \in J} X _ {i} \subset \bigcup_ {i \in I} X _ {i}, \quad \text { and } (\text { if } J \neq \emptyset) \quad \bigcap_ {i \in J} X _ {i} \supset \bigcap_ {i \in I} X _ {i}.
\]

PROPOSITION 2. Let \((\mathbf{X}_{\iota})_{\iota \in \mathbf{I}}\) be a family of sets whose index set I is the union of a family \((\mathrm{J}_{\lambda})_{\lambda \in \mathbf{L}}\) of sets. Then

\[
\bigcup_ {i \in I} X _ {i} = \bigcup_ {\lambda \in L} \left(\bigcup_ {i \in J _ {\lambda}} X _ {i}\right),
\]

and (if \(\mathbf{L} \neq \phi\) and \(J_{\lambda} \neq \phi\) for each \(\lambda \in \mathbf{L}\) )

\[
\bigcap_ {i \in I} X _ {i} = \bigcap_ {\lambda \in L} \left(\bigcap_ {i \in J _ {\lambda}} X _ {i}\right)
\]

(``associativity'' of union and intersection).

Let \(x\) be an element of \(\bigcup_{i \in I} X_i\) . There exists an index \(i \in I\) such that \(x \in X_i\) . Since \(I\) is the union of the family \((J_\lambda)_{\lambda \in L}\) , there exists an index \(\lambda \in L\) such that \(i \in J_\lambda\) , whence \(x \in \bigcup X_i\) , and consequently

\[
x \in \bigcup_ {\lambda \in L} \left(\bigcup_ {i \in J _ {\lambda}} X _ {i}\right).
\]

Conversely, let \(x\) be an element of this set. There exists an index \(\lambda \in L\) such that \(x \in \bigcup_{\iota \in J_{\lambda}} X_{\iota}\) , hence there exists an index \(\iota \in J_{\lambda}\) (and therefore \(\iota \in I\) ) such that \(x \in X_{\iota}\) ; it follows that

\[
x \in \bigcup_ {i \in I} X _ {i}.
\]

Now suppose that \(L \neq \emptyset\) and \(J_{\lambda} \neq \emptyset\) for each \(\lambda \in L\) . Then \(I \neq \emptyset\) . Let x be an element of \(\bigcap_{i \in I} X_{i}\) . If \(\lambda \in L\) , we have \(x \in X_{i}\) for each \(i \in J_{\lambda}\) (since \(J_{\lambda} \subset I\) ), whence \(x \in \bigcap_{i \in J_{\lambda}} X_{i}\) . Since the last inclusion is true for all \(\lambda \in L\) , x belongs to \(\bigcap_{\lambda \in L} \left( \bigcap_{i \in J_{\lambda}} X_{i} \right)\) . Conversely, let x be an element of this latter set, and let i be any element of I. There exists \(\lambda \in L\) such that \(i \in J_{\lambda}\) ; since \(x \in \bigcap_{i \in J_{\lambda}} X_{i}\) , we have \(x \in X_{i}\) , which is true for all \(i \in I\) . Hence \(x \in \bigcap_{i \in I} X_{i}\) . This completes the proof.

For families of subsets of a given set the second part of Proposition 2 remains valid without restriction on L and \(J_{\lambda}\) .

